\subsection{Approximation-guaranteed greedy construction}
\label{subsec:aor-weighted-greedy}

The identity-closure reduction also permits a fast constructive method with a
classical performance guarantee.  Fix all mandatory strategies, including the
inactive strategy, and remove the tagged requirements they cover.  Let
$U^0$ be the residual requirement universe, let $W_0\geq 0$ be the fixed
mandatory burden, and let each residual strategy $s$ have positive exact
weight $w_s$ and residual coverage $R_s\subseteq U^0$.  For nonempty $U^0$,
write
\[
 d=\max_s |R_s|,
 \qquad
 H(d)=\sum_{i=1}^d\frac1i.
\]

The weighted rule is the standard set-cover greedy algorithm of
\citet{Chvatal1979}, applied only after the semantic reduction:

\begin{enumerate}
\item initialize the selected set with every mandatory strategy and set
      $U\leftarrow U^0$;
\item while $U\neq\varnothing$, choose a residual strategy minimizing
      \[
        \frac{w_s}{|R_s\cap U|}
      \]
      among strategies with nonzero new coverage, breaking an exact tie by
      the lowest original strategy index in canonical identifier order;
\item add the chosen strategy and replace
      $U\leftarrow U\setminus R_s$;
\item optionally examine the chosen nonmandatory strategies in reverse
      selection order and delete one exactly when complete tagged coverage
      remains.
\end{enumerate}

All scores and comparisons in the implementation use
\texttt{Rational\{BigInt\}}.  A separate cardinality variant replaces
$w_s$ in the score by one, equivalently choosing the largest number of newly
covered requirements.  That variant has the burden guarantee below only when
all residual optional weights are equal; with unequal weights, no
approximation guarantee for the actual retention burden is asserted.

\begin{theorem}[Transferred weighted-greedy guarantee]
\label{thm:aor-weighted-greedy}
Suppose the source library is finite, closure is identity closure, every
residual tagged requirement has a carrier, mandatory weights are nonnegative,
and every nonmandatory weight is positive.  If $U^0\neq\varnothing$, the
weighted greedy construction returns an innovation-safe source sublibrary of
burden $W_G$ satisfying
\[
 W_G\leq H(d)\,\operatorname{OPT},
\]
where $\operatorname{OPT}$ is the minimum total burden of any safe source
sublibrary and $d=\max_s|R_s|$.  Deterministic reverse deletion returns burden
$W_{GR}\leq W_G$ and hence satisfies the same bound.  If $U^0$ is empty, the
mandatory library is optimal and the factor is one.
\end{theorem}

\begin{proof}
First suppose $U^0$ is nonempty.  The residual instance is an ordinary
positive-cost weighted set-cover instance.  Chv\'atal's theorem applies to
the rule minimizing cost per newly covered element and gives
\[
 G\leq H(d)\operatorname{OPT}_R,
\]
where $G$ and $\operatorname{OPT}_R$ are, respectively, the greedy and
optimal residual costs.  The theorem permits any choice among exact
minimum-ratio ties, so the canonical deterministic tie rule does not change
the bound.

Every admissible source selection contains the mandatory strategies, and
mandatory-only preprocessing removes no optional strategy.  Thus
$W_G=W_0+G$ and
$\operatorname{OPT}=W_0+\operatorname{OPT}_R$.  Since $W_0\geq0$ and
$H(d)\geq1$,
\[
 W_G=W_0+G
 \leq W_0+H(d)\operatorname{OPT}_R
 \leq H(d)(W_0+\operatorname{OPT}_R)
 =H(d)\operatorname{OPT}.
\]
The tagged-cover equivalence then converts residual complete coverage plus
mandatory retention into exact preservation of the source operating frontier
and source identity closure.

Reverse deletion accepts a removal only after complete tagged coverage is
rechecked.  It never removes a mandatory strategy, and every removed optional
strategy has positive weight.  Therefore each accepted deletion strictly
lowers burden, every rejected deletion leaves burden unchanged, and
$W_{GR}\leq W_G$.  Moreover, one reverse pass is enough for
inclusion-wise irreducibility: if a strategy cannot be removed at its
examination, later deletions cannot restore any missing coverage.  The same
$H(d)$ bound follows immediately.  If $U^0$ is empty, mandatory retention is
already feasible and every optional positive-weight strategy is unnecessary,
so the mandatory library is the exact optimum.
\end{proof}

\paragraph{A primary-source near-tight fixture.}
Chv\'atal's example has $m$ requirements, singleton set $\{j\}$ of cost
$1/j$ for each $j=1,\ldots,m$, and a set covering the entire universe with
cost $1+\epsilon$, where $\epsilon>0$ and
$1+\epsilon<H(m)$.  Greedy selects the singletons in order
$m,m-1,\ldots,1$: when $q$ requirements remain, singleton $q$ has score
$1/q$, whereas the universe set has score $(1+\epsilon)/q>1/q$.  Its burden
is $H(m)$, while the universe set alone is optimal with burden
$1+\epsilon$.  The ratio $H(m)/(1+\epsilon)$ tends to $H(m)$ as
$\epsilon\downarrow0$.  The journal fixture realizes each requirement as an
identity-closure module tag and lets the mandatory zero-weight inactive
strategy cover the sole zero-frontier tag.  Reverse deletion removes no
singleton from the greedy endpoint.

\begin{example}[Shared three-row greedy arithmetic]
\label{ex:aor-greedy-shared-formal-fixture}
For \(m=3\) and \(\epsilon=1/12\), the singleton burden is
\(1+1/2+1/3=11/6\), the bundle burden is \(1+1/12=13/12\), and the exact
ratio is \(22/13\).  The Julia--Lean shared fixture checks only this arithmetic;
it does not formalize Chv\'atal's theorem or the transfer to safe compression.
\end{example}

\paragraph{Limitations and evidence boundary.}
The guarantee is inherited set-cover theory; the contribution here is its
proved transfer through the source-relative frontier--closure representation
and exact postchecked implementation.  The theorem does not apply to
arbitrary nonidentity closure, because independent module rows need not
represent closure complementarities.  It does not give the backward
heaviest-safe-first deletion heuristic an approximation guarantee, and it
does not make cardinality greedy a weighted-burden approximation when weights
differ.  The returned certificate is exact finite computation: it rechecks
mandatory retention, tagged coverage, the original frontier, the original
identity closure, and burden.  It is neither exhaustive optimality evidence,
mixed-integer solver evidence, nor Lean kernel verification.
